\section{Conclusion}

In this work, we have constructed a theory capable of formally describing classical Go from elementary notions—\emph{where, who, what, and how}—supported by set theory. Throughout this construction, the rules of the game emerge naturally, and intuitive concepts acquire a precise and tractable expression.

However, the principal value of this framework is not limited to formalizing Go. As explored in Part~\ref{parte3}, the developed structure is sufficiently general to encompass a broader family of games based on the idea of surrounding and capturing. This suggests the existence of a more abstract mathematical object, which we designate as \textsc{Surrounding Games} (\(\mathsf{SG}\)), of which Go is merely a particular case.

Several paths remain open: identifying which properties are exclusive to Go and which belong to the entire family, formalizing this family and seeing how far we can take it, constructing infinite Go, or attempting to connect this theory with combinatorial game theory.

\vfill

I've realize that most of my work render the 60/70 papers, when I look up for sources, those were the ones that, in my opinion it had the ideas and concepts I thought will be usefull or necesary. There many more papers that at first glance seems to be for game theory but they do care about the game as a model. I'll read those and try to make this work a coherent, simple (I hope) and pretty book of the game.